\section{Core Concepts and \acr{LLVM} Syntax}

This section introduces the fundamental concepts and syntax elements that form the foundation of \acr{LLVM} \acr{IR} programming. Some concepts are unique to \acr{LLVM}, while others derive from established compiler theory; we explain the rules and conventions for their usage within \acr{LLVM} \acr{IR}.

\subsection{Static Single Assignment (\acr{SSA}) Form}

\acr{SSA} is a property where each variable is assigned exactly once and every use refers to a specific definition. This constraint makes data flow explicit and enables powerful compiler optimizations by simplifying the analysis of how values propagate through the program.

\textbf{Non-\acr{SSA} code:}
\begin{lstlisting}[style=ir]
x = 1
x = x + 2  // x is assigned twice (not SSA)
\end{lstlisting}

\textbf{\acr{SSA} code:}
\begin{lstlisting}[style=ir]
%x1 = 1
%x2 = add i32 %x1, 2  // Each value has unique name
\end{lstlisting}

The benefits of \acr{SSA} form include:
\begin{itemize}[noitemsep]
  \item \textbf{Constant propagation}: The value of \texttt{\%x2} is provably 3, enabling compile-time evaluation
  \item \textbf{Dead code elimination}: Unused definitions can be safely removed without complex liveness analysis
  \item \textbf{Register allocation}: Clear lifetime of each value simplifies mapping virtual registers to physical registers
  \item \textbf{Simplified dependency tracking}: Each definition has a unique name, making def-use chains trivial to construct
\end{itemize}

\subsection{Modules and Target Specification}

\acr{LLVM} \acr{IR} is organized hierarchically, with the \textbf{module} serving as the top-level container. A module corresponds to a single compilation unit and contains functions, global variables, type definitions, and metadata.

By convention, \acr{LLVM} \acr{IR} files begin with comments identifying the module and source:

\begin{lstlisting}[style=ir]
; ModuleID = 'my_module.ll'
; source_filename = "my_source.c"
\end{lstlisting}

These comments are automatically generated by compilers and help identify the module's origin during debugging. The \texttt{@euriklis/llvm-ir} library provides dedicated methods to generate this metadata programmatically, ensuring proper module identification in generated code.

Every module must specify its target architecture through two declarations:

\begin{lstlisting}[style=ir]
target triple = "x86_64-unknown-linux-gnu"
target datalayout = "e-m:e-p270:32:32-i64:64-f80:128-n8:16:32:64-S128"
\end{lstlisting}

The \textbf{target triple} follows the format \texttt{architecture-vendor-os-environment} and identifies the platform:

\begin{itemize}[noitemsep]
  \item \texttt{x86\_64-unknown-linux-gnu} --- Linux on x86-64
  \item \texttt{\acr{x86}\_64-apple-darwin} --- macOS on Intel
  \item \texttt{aarch64-unknown-linux-gnu} --- Linux on \acr{ARM} (\acr{AWS} Graviton, Raspberry Pi)
  \item \texttt{aarch64-apple-darwin} --- macOS on Apple Silicon
  \item \texttt{i686-unknown-linux-gnu} --- Linux on x86 32-bit
  \item \texttt{riscv64-unknown-linux-gnu} --- Linux on \acr{RISC-V} 64-bit
  \item \texttt{powerpc64le-unknown-linux-gnu} --- Linux on PowerPC 64-bit little-endian
  \item \texttt{nvptx64-nvidia-cuda} --- NVIDIA \acr{GPU}s
  \item \texttt{amdgcn-amd-amdhsa} --- \acr{AMD} \acr{GPU}s (\acr{ROCm})
  \item \texttt{wasm32-unknown-unknown} --- \acr{WASM}
  \item \texttt{spirv64-unknown-unknown} --- \acr{SPIR-V} for Vulkan/\acr{OpenCL}
\end{itemize}

To view the complete list of supported target triples on your system, run \texttt{llc --version}, which displays all registered targets. For detailed information about specific targets, use \texttt{llc -march=<arch> -mattr=help} to see available CPU features and subtargets.

\textbf{Example (LLVM 18.1.3):}

\begin{lstlisting}[style=ir]
$ llc --version
Ubuntu LLVM version 18.1.3
  Optimized build.
  Default target: x86_64-pc-linux-gnu
  Host CPU: znver2

  Registered Targets:
    aarch64     - AArch64 (little endian)
    aarch64_32  - AArch64 (little endian ILP32)
    aarch64_be  - AArch64 (big endian)
    amdgcn      - AMD GCN GPUs
    arm         - ARM
    arm64       - ARM64 (little endian)
    arm64_32    - ARM64 (little endian ILP32)
    armeb       - ARM (big endian)
    avr         - Atmel AVR Microcontroller
    bpf         - BPF (host endian)
    bpfeb       - BPF (big endian)
    bpfel       - BPF (little endian)
    hexagon     - Hexagon
    lanai       - Lanai
    loongarch32 - 32-bit LoongArch
    loongarch64 - 64-bit LoongArch
    m68k        - Motorola 68000 family
    mips        - MIPS (32-bit big endian)
    mips64      - MIPS (64-bit big endian)
    mips64el    - MIPS (64-bit little endian)
    mipsel      - MIPS (32-bit little endian)
    msp430      - MSP430 [experimental]
    nvptx       - NVIDIA PTX 32-bit
    nvptx64     - NVIDIA PTX 64-bit
    ppc32       - PowerPC 32
    ppc32le     - PowerPC 32 LE
    ppc64       - PowerPC 64
    ppc64le     - PowerPC 64 LE
    r600        - AMD GPUs HD2XXX-HD6XXX
    riscv32     - 32-bit RISC-V
    riscv64     - 64-bit RISC-V
    sparc       - Sparc
    sparcel     - Sparc LE
    sparcv9     - Sparc V9
    systemz     - SystemZ
    thumb       - Thumb
    thumbeb     - Thumb (big endian)
    ve          - VE
    wasm32      - WebAssembly 32-bit
    wasm64      - WebAssembly 64-bit
    x86         - 32-bit X86: Pentium-Pro and above
    x86-64      - 64-bit X86: EM64T and AMD64
    xcore       - XCore
    xtensa      - Xtensa 32
\end{lstlisting}

\textbf{Example of target-specific features (NVIDIA PTX):}

\begin{lstlisting}[style=ir]
$ llc -march=nvptx64 -mattr=help
Available CPUs for this target:

  sm_20  - Select the sm_20 processor.
  sm_21  - Select the sm_21 processor.
  sm_30  - Select the sm_30 processor.
  sm_32  - Select the sm_32 processor.
  sm_35  - Select the sm_35 processor.
  sm_37  - Select the sm_37 processor.
  sm_50  - Select the sm_50 processor.
  sm_52  - Select the sm_52 processor.
  sm_53  - Select the sm_53 processor.
  sm_60  - Select the sm_60 processor.
  sm_61  - Select the sm_61 processor.
  sm_62  - Select the sm_62 processor.
  sm_70  - Select the sm_70 processor.
  sm_72  - Select the sm_72 processor.
  sm_75  - Select the sm_75 processor.
  sm_80  - Select the sm_80 processor.
  sm_86  - Select the sm_86 processor.
  sm_87  - Select the sm_87 processor.
  sm_89  - Select the sm_89 processor.
  sm_90  - Select the sm_90 processor.
  sm_90a - Select the sm_90a processor.

Available features for this target:

  ptx32  - Use PTX version 32.
  ptx40  - Use PTX version 40.
  ptx41  - Use PTX version 41.
  ptx42  - Use PTX version 42.
  ptx43  - Use PTX version 43.
  ptx50  - Use PTX version 50.
  ptx60  - Use PTX version 60.
  ptx61  - Use PTX version 61.
  ptx63  - Use PTX version 63.
  ptx64  - Use PTX version 64.
  ptx65  - Use PTX version 65.
  ptx70  - Use PTX version 70.
  ptx71  - Use PTX version 71.
  ptx72  - Use PTX version 72.
  ptx73  - Use PTX version 73.
  ptx74  - Use PTX version 74.
  ptx75  - Use PTX version 75.
  ptx76  - Use PTX version 76.
  ptx77  - Use PTX version 77.
  ptx78  - Use PTX version 78.
  ptx80  - Use PTX version 80.
  ptx81  - Use PTX version 81.
  ptx82  - Use PTX version 82.
  ptx83  - Use PTX version 83.
  sm_20  - Target SM 20.
  sm_21  - Target SM 21.
  sm_30  - Target SM 30.
  sm_32  - Target SM 32.
  sm_35  - Target SM 35.
  sm_37  - Target SM 37.
  sm_50  - Target SM 50.
  sm_52  - Target SM 52.
  sm_53  - Target SM 53.
  sm_60  - Target SM 60.
  sm_61  - Target SM 61.
  sm_62  - Target SM 62.
  sm_70  - Target SM 70.
  sm_72  - Target SM 72.
  sm_75  - Target SM 75.
  sm_80  - Target SM 80.
  sm_86  - Target SM 86.
  sm_87  - Target SM 87.
  sm_89  - Target SM 89.
  sm_90  - Target SM 90.
  sm_90a - Target SM 90a.

Use +feature to enable a feature, or -feature to disable it.
For example, llc -mcpu=mycpu -mattr=+feature1,-feature2
\end{lstlisting}

The \textbf{data layout string} is a compact specification that encodes essential low-level architectural properties required for correct code generation~\cite{llvmlangref}. This string informs the \acr{LLVM} optimizer and backend about target-specific characteristics such as endianness, pointer sizes, type alignments, and memory layout conventions~\cite{lopesauler2014}. Without accurate data layout information, the compiler cannot generate correct memory access patterns, properly align data structures, or make sound optimization decisions that depend on type sizes and alignment constraints~\cite{llvmessentials}.

Understanding the data layout string reveals how the same \acr{IR} code can behave differently across platforms. Consider a simple operation like storing a 32-bit integer: the compiler must know \textbf{endianness} (byte ordering in memory), \textbf{alignment} (valid memory addresses for types), and \textbf{pointer size} (32-bit vs. 64-bit addressing). These architectural differences are abstracted away in high-level languages but become critical when generating machine code. When \acr{LLVM} compiles your \acr{IR}, it uses the data layout to:

\begin{enumerate}[noitemsep]
  \item \textbf{Generate correct memory instructions}: Determine how to load/store multi-byte values (e.g., reading a 64-bit integer requires knowing if bytes are stored least-significant-first or most-significant-first)
  \item \textbf{Calculate struct padding}: Insert padding bytes between struct fields to meet alignment requirements, ensuring efficient memory access
  \item \textbf{Optimize memory access patterns}: Use \acr{SIMD} instructions when data is properly aligned, or fall back to slower unaligned access when necessary
  \item \textbf{Implement pointer arithmetic}: Correctly compute memory offsets for array indexing and structure field access based on target-specific type sizes
\end{enumerate}

Think of the data layout as a contract between your \acr{IR} code and the target hardware---it ensures the compiled binary correctly interacts with physical memory and CPU registers.

\subsubsection*{Data Layout String Format}

The data layout specification consists of multiple components separated by the minus sign (\texttt{-}), where each component begins with a letter identifying the specification type, followed by parameters~\cite{llvmlangref}. A typical \acr{x86}-64 data layout string is:

\begin{lstlisting}[style=ir]
target datalayout = "e-m:e-p270:32:32-p271:32:32-p272:64:64-i64:64-f80:128-n8:16:32:64-S128"
\end{lstlisting}

\subsubsection*{Component Breakdown}

\textbf{Endianness (\texttt{e} or \texttt{E}):} Endianness defines the order in which bytes of a multi-byte value are stored in memory. When you store a number larger than one byte (like a 32-bit integer), the CPU must decide: should the first byte in memory be the most significant or least significant? Different CPU architectures make different choices.

The first character in the data layout specifies byte ordering~\cite{llvmlangref}:
\begin{itemize}[noitemsep]
  \item \texttt{e} --- Little-endian: least significant byte at lowest memory address (Intel/\acr{AMD} \acr{x86}, \acr{ARM} in little-endian mode)
  \item \texttt{E} --- Big-endian: most significant byte at lowest memory address (legacy PowerPC, SPARC, some network protocols)
\end{itemize}

\textbf{Example:} Consider the 32-bit integer \texttt{0x12345678}. This number occupies 4 bytes in memory:

\begin{lstlisting}[style=ir]
Address:  0x1000  0x1001  0x1002  0x1003
Little:     0x78    0x56    0x34    0x12   (least significant first)
Big:        0x12    0x34    0x56    0x78   (most significant first)
\end{lstlisting}

Why it matters: When reading data from disk, network, or interfacing with hardware, the endianness must match. Reading a little-endian value as big-endian would interpret \texttt{0x12345678} as \texttt{0x78563412}---completely wrong!

\textbf{Name Mangling (\texttt{m:<style>}):} Specifies how symbol names are mangled for the linker~\cite{llvmlangref}:
\begin{itemize}[noitemsep]
  \item \texttt{m:e} --- ELF mangling (Linux, BSD, most Unix systems)
  \item \texttt{m:o} --- Mach-O mangling (macOS, iOS)
  \item \texttt{m:m} --- Mips mangling
  \item \texttt{m:w} --- Windows COFF mangling
  \item \texttt{m:x} --- Windows COFF \acr{x86} mangling (different from \texttt{m:w} for compatibility)
\end{itemize}

Symbol mangling affects how function and global variable names are encoded in object files, determining linker compatibility and calling convention details.

\textbf{Pointer Layout (\texttt{p[n]:<size>:<abi>:<pref>}):} Specifies pointer characteristics for address space \texttt{n}~\cite{llvmlangref}. If \texttt{n} is omitted, address space 0 (default) is assumed:
\begin{itemize}[noitemsep]
  \item \texttt{<size>} --- Pointer size in bits
  \item \texttt{<abi>} --- \acr{ABI}-required alignment in bits (minimum guaranteed by calling conventions)
  \item \texttt{<pref>} --- Preferred alignment in bits for performance (optional; if omitted, equals \texttt{<abi>})
\end{itemize}

Examples from the \acr{x86}-64 data layout:
\begin{itemize}[noitemsep]
  \item \texttt{p270:32:32} --- Address space 270: 32-bit pointers with 32-bit alignment (used for \texttt{\_\_ptr32} Microsoft extension)
  \item \texttt{p271:32:32} --- Address space 271: 32-bit zero-extended pointers (used for \texttt{\_\_uptr})
  \item \texttt{p272:64:64} --- Address space 272: 64-bit pointers with 64-bit alignment (used for \texttt{\_\_ptr64})
  \item Default (\texttt{p:64:64:64}) --- Address space 0: 64-bit pointers with 64-bit alignment (standard pointers)
\end{itemize}

Multiple address spaces enable heterogeneous programming where different pointer types have different sizes (e.g., \acr{GPU} programming with separate host and device address spaces)~\cite{lopesauler2014}.

\textbf{Integer Alignment (\texttt{i<size>:<abi>:<pref>}):} Alignment specifies which memory addresses are valid for storing a value. A value is \textbf{aligned} if its memory address is divisible by its size. For example, a 32-bit (4-byte) integer should be stored at addresses divisible by 4 (like 0x1000, 0x1004, 0x1008), not at arbitrary addresses like 0x1001.

The alignment specification format is \texttt{i<size>:<abi>:<pref>}~\cite{llvmlangref}:

\begin{lstlisting}[style=ir]
i64:64      ; i64 has 64-bit (8-byte) ABI and preferred alignment
i32:32      ; i32 has 32-bit (4-byte) alignment
i8:8        ; i8 has 8-bit (1-byte) alignment
\end{lstlisting}

\textbf{Example:} Storing an \texttt{i32} (4-byte integer) in memory:

\begin{lstlisting}[style=ir]
Valid addresses (aligned):   0x1000, 0x1004, 0x1008, 0x100C
Invalid addresses (misaligned): 0x1001, 0x1002, 0x1003, 0x1005
\end{lstlisting}

Why alignment matters:
\begin{itemize}[noitemsep]
  \item \textbf{Performance}: CPUs read memory in aligned chunks (cache lines). An aligned 64-bit read takes one memory transaction; a misaligned read may require two transactions, doubling access time
  \item \textbf{Correctness}: Some CPUs (ARM, SPARC) raise hardware exceptions or crashes on misaligned access
  \item \textbf{Atomic operations}: CPU atomic instructions (compare-and-swap, etc.) require natural alignment; misaligned atomics may silently produce incorrect results
\end{itemize}

\textbf{Floating-Point Alignment (\texttt{f<size>:<abi>:<pref>}):} Specifies alignment for floating-point types~\cite{llvmlangref}:

\begin{lstlisting}[style=ir]
f80:128     ; x86 80-bit extended precision (stored in 128-bit slots)
f64:64      ; double (64-bit float)
f32:32      ; float (32-bit float)
\end{lstlisting}

The \texttt{f80:128} entry indicates that while \acr{x86} extended precision floats are 80 bits, they are stored in 128-bit-aligned memory slots for efficiency.

\textbf{Vector Alignment (\texttt{v<size>:<abi>:<pref>}):} Specifies alignment for vector types used in \acr{SIMD} operations~\cite{llvmlangref}:

\begin{lstlisting}[style=ir]
v128:128    ; 128-bit vectors (SSE on x86, NEON on ARM)
v256:256    ; 256-bit vectors (AVX)
v512:512    ; 512-bit vectors (AVX-512)
\end{lstlisting}

Proper vector alignment is mandatory for \acr{SIMD} instruction efficiency. Misaligned vector loads/stores can cause significant performance penalties or instruction faults.

\textbf{Aggregate Alignment (\texttt{a:<abi>:<pref>}):} Specifies alignment for aggregate types (structs, arrays)~\cite{llvmlangref}:

\begin{lstlisting}[style=ir]
a:0:64      ; Aggregates have 0-bit ABI alignment, 64-bit preferred alignment
\end{lstlisting}

This allows the compiler to optimize struct layout while maintaining \acr{ABI} compatibility.

\textbf{Native Integer Widths (\texttt{n<size1>:<size2>:...}):} Specifies integer widths that the target CPU can efficiently process in a single instruction~\cite{llvmlangref}:

\begin{lstlisting}[style=ir]
n8:16:32:64 ; x86-64 efficiently supports 8, 16, 32, and 64-bit integers
n32:64      ; PowerPC 64 efficiently supports 32 and 64-bit integers
\end{lstlisting}

The optimizer uses this information for type legalization decisions. Operations on \textbf{\textit{non-native widths}} (e.g., \texttt{i7}, \texttt{i128} on \acr{x86}-64) are decomposed into multiple native-width operations.

\textbf{Stack Alignment (\texttt{S<size>}):} Specifies the natural stack alignment in bits~\cite{llvmlangref}:

\begin{lstlisting}[style=ir]
S128        ; Stack maintains 128-bit (16-byte) alignment
\end{lstlisting}

Stack alignment affects:
\begin{itemize}[noitemsep]
  \item Function prologues/epilogues (stack frame setup)
  \item \acr{SIMD} variable allocation (local vectors must be aligned)
  \item Calling conventions (some \acr{ABI}s require specific stack alignment)
\end{itemize}

On \acr{x86}-64, 16-byte stack alignment enables efficient \acr{SSE} operations on stack-allocated vectors~\cite{lopesauler2014}.

The data layout string is essential for three fundamental reasons~\cite{llvmessentials}: \textbf{(1) Correctness}---wrong alignment can cause crashes on strict-alignment architectures or incorrect atomic operations; \textbf{(2) Performance}---proper alignment enables efficient memory access, \acr{SIMD} vectorization, and cache utilization; and \textbf{(3) \acr{ABI} compliance}---data layout must match the platform \acr{ABI} to ensure interoperability with system libraries and other compiled code. Modern \acr{LLVM} automatically sets the data layout when you specify the target triple, but understanding the format is crucial for cross-compilation, \acr{GPU} programming, and custom target development~\cite{lopesauler2014}.

\subsection{Global and Local Variables}

\acr{LLVM} distinguishes between global variables (module-level storage) and local variables (function-level virtual registers). In \acr{LLVM} terminology, these constructs are called \textbf{identifiers}. Global identifiers (functions, global variables) begin with the \texttt{@} character, while local identifiers (virtual registers, local names) begin with the \texttt{\%} character. This prefix-based naming scheme serves a crucial design purpose: it prevents conflicts with reserved words (such as opcodes like \texttt{add}, \texttt{mul}, or type names like \texttt{i32}, \texttt{void}). Since reserved words in \acr{LLVM} never start with \texttt{@} or \texttt{\%}, identifiers can never clash with them. This design provides flexibility\,---\,\acr{LLVM} can expand its set of reserved words in future versions without breaking existing code, as no valid identifier can conflict with a reserved word.

Identifiers come in two forms: \textbf{named} (e.g., \texttt{\%foo}, \texttt{@DivisionByZero}) and \textbf{unnamed} (e.g., \texttt{\%12}, \texttt{@2}). Named identifiers follow the regular expression pattern \texttt{[-a-zA-Z\$.\allowbreak\_][-a-zA-Z\$.\allowbreak\_0-9]*}. Unnamed identifiers allow compilers to quickly generate temporary values without symbol table lookups.

\textbf{Global variables} are declared outside functions and persist for the program's lifetime:

\begin{lstlisting}[style=ir]
@global_counter = global i32 0
@constant_pi = constant float 3.14159
@external_array = external global [100 x i32]
\end{lstlisting}

Global variables use the \texttt{@} prefix and have linkage specifiers:
\begin{itemize}[noitemsep]
  \item \texttt{global} --- Mutable global variable
  \item \texttt{constant} --- Immutable global constant
  \item \texttt{external} --- Declared elsewhere (no initializer)
  \item \texttt{internal} --- Private to the module
\end{itemize}

\textbf{Local variables} (virtual registers) use the \texttt{\%} prefix and exist only within a function:

\begin{lstlisting}[style=ir]
define i32 @compute(i32 %x) {
entry:
  %temp1 = add i32 %x, 5
  %temp2 = mul i32 %temp1, 2
  ret i32 %temp2
}
\end{lstlisting}

In \acr{SSA} form, each local variable is assigned exactly once. For mutable variables (stack allocations), \acr{LLVM} uses memory operations:

\begin{lstlisting}[style=ir]
define void @mutable_example() {
entry:
  %ptr = alloca i32              ; Allocate stack space
  store i32 42, ptr %ptr         ; Write to memory
  %val = load i32, ptr %ptr      ; Read from memory
  ret void
}
\end{lstlisting}

\subsubsection*{SSA Values vs. Memory Locations}

Understanding the distinction between \acr{SSA} values and memory locations is crucial:

\textbf{SSA values} (virtual registers with \texttt{\%} prefix) can only be assigned \textbf{once}:

\begin{lstlisting}[style=ir]
%x = add i32 1, 2      ; %x = 3
%x = mul i32 4, 5      ; ERROR: Cannot reassign %x (violates SSA)
\end{lstlisting}

Each new value requires a new variable name:

\begin{lstlisting}[style=ir]
%x = add i32 1, 2      ; %x = 3
%y = mul i32 %x, 5     ; %y = 15 (new variable name)
%z = add i32 %y, 1     ; %z = 16 (new variable name)
\end{lstlisting}

\textbf{Memory locations} (globals and \texttt{alloca}) can be stored to \textbf{multiple times}:

\begin{lstlisting}[style=ir]
@global_var = global i32 0

define void @modify_memory() {
  %local_ptr = alloca i32

  ; Multiple stores to global memory
  store i32 10, ptr @global_var
  store i32 20, ptr @global_var    ; OK: Overwrites previous value
  store i32 30, ptr @global_var    ; OK: Overwrites again

  ; Multiple stores to local memory
  store i32 100, ptr %local_ptr
  store i32 200, ptr %local_ptr    ; OK: Overwrites previous value

  ret void
}
\end{lstlisting}

\textbf{Key distinction}: The pointer variable (\texttt{\%local\_ptr}) is assigned once (the \texttt{alloca} result), but the \textbf{memory location} it points to can be modified multiple times. Each \texttt{load} creates a new \acr{SSA} value:

\begin{lstlisting}[style=ir]
%ptr = alloca i32              ; %ptr assigned once (address)
store i32 10, ptr %ptr
%val1 = load i32, ptr %ptr     ; %val1 = 10 (new SSA value)
store i32 20, ptr %ptr         ; Modify memory
%val2 = load i32, ptr %ptr     ; %val2 = 20 (new SSA value)
\end{lstlisting}

This pattern allows mutable variables while maintaining \acr{SSA} form: the pointer names follow \acr{SSA} rules, but the memory they point to can change.

\subsection{Functions}

Functions are the primary units of executable code in \acr{LLVM} \acr{IR}. A function declaration consists of:

\begin{lstlisting}[style=ir]
define <return_type> @function_name(<param_type> %param1, ...) {
  ; function body
}
\end{lstlisting}

The observant reader has probably noticed that the naming convention for functions in \acr{LLVM} matches that of global variables---both use the \texttt{@} prefix. This similarity is not coincidental: \textbf{functions are global values} in \acr{LLVM} \acr{IR}. Like global variables, functions are:

\begin{itemize}[noitemsep]
  \item Declared at \textbf{module scope} (not inside other functions)
  \item Visible across the entire module
  \item Have \textbf{addresses} that can be referenced and stored
  \item Support \textbf{linkage types} controlling visibility and linking behavior
\end{itemize}

Functions can be treated as values and stored in function pointers:

\begin{lstlisting}[style=ir]
define i32 @my_function(i32 %x) {
  %result = add i32 %x, 1
  ret i32 %result
}

@function_ptr = global ptr @my_function  ; Store function address

define void @call_via_pointer() {
  %fptr = load ptr, ptr @function_ptr
  %result = call i32 %fptr(i32 42)       ; Call through pointer
  ret void
}
\end{lstlisting}

In \acr{LLVM}'s type hierarchy, both \texttt{Function} and \texttt{GlobalVariable} inherit from \texttt{GlobalValue}, making them module-level entities with global visibility.

\textbf{Linkage types} control visibility and linking behavior:

\begin{itemize}[noitemsep]
  \item \texttt{define} --- Default linkage (internal to module unless exported)
  \item \texttt{define external} --- Exported symbol, visible to other modules
  \item \texttt{define internal} --- Private to the module (similar to C \texttt{static})
  \item \texttt{define weak} --- Can be overridden by other definitions
  \item \texttt{declare} --- External function (declaration only, no body)
\end{itemize}

\textbf{Function attributes} modify compilation behavior:

\begin{lstlisting}[style=ir]
define i32 @pure_function(i32 %x) nounwind readnone {
  %result = mul i32 %x, %x
  ret i32 %result
}
\end{lstlisting}

Common attributes:
\begin{itemize}[noitemsep]
  \item \texttt{nounwind} --- Function never throws exceptions
  \item \texttt{readnone} --- Pure function, no memory access
  \item \texttt{readonly} --- Function only reads memory
  \item \texttt{alwaysinline} --- Force inlining
  \item \texttt{noinline} --- Prevent inlining
\end{itemize}

\subsection{Type System}

\acr{LLVM} \acr{IR} is strongly typed. Every value has an explicit type, and type mismatches are compile-time errors. The type system includes primitive types, aggregate types, and pointer types.

\textbf{Primitive types:}

\begin{itemize}[noitemsep]
  \item \textbf{Integer types}: \texttt{i1} (boolean), \texttt{i8}, \texttt{i16}, \texttt{i32}, \texttt{i64}, \texttt{i128} (arbitrary-width integers like \texttt{i7} or \texttt{i256} are also valid)
  \item \textbf{Floating-point types}: \texttt{half} (16-bit), \texttt{float} (32-bit), \texttt{double} (64-bit), \texttt{fp128} (quad-precision)
  \item \textbf{Void type}: \texttt{void} (used only for function return types, not for values)
\end{itemize}

\textbf{Aggregate types:}

\begin{lstlisting}[style=ir]
; Array: [length x element_type]
[10 x i32]                    ; Array of 10 i32 values
[3 x float]                   ; Array of 3 floats

; Structure: { type1, type2, ... }
{ i32, float, ptr }           ; Unpacked struct
<{ i8, i32 }>                ; Packed struct (no padding)

; Vector: <length x element_type>
<4 x float>                   ; SIMD vector (SSE, NEON)
<16 x i8>                     ; 128-bit packed integer vector
\end{lstlisting}

\textbf{Pointer type}: \texttt{ptr} (opaque pointer, introduced in \acr{LLVM} 15):

\begin{lstlisting}[style=ir]
%heap_ptr = call ptr @malloc(i64 64)
store i32 42, ptr %heap_ptr
%value = load i32, ptr %heap_ptr
\end{lstlisting}

Older \acr{LLVM} versions used typed pointers (\texttt{i32*}, \texttt{float*}), but modern \acr{LLVM} uses opaque \texttt{ptr} to simplify the type system and improve optimization.

\textbf{Function types} specify signatures:

\begin{lstlisting}[style=ir]
; Function type: return_type (param_types)
i32 (i32, i32)                ; Takes two i32, returns i32
void (ptr, i64)               ; Takes pointer and i64, returns void
\end{lstlisting}

\subsection{Basic Blocks and Labels}

\acr{LLVM} organizes function bodies into \textbf{basic blocks}---straight-line sequences of instructions with no branches except at the end. This structure is fundamental to control flow analysis and optimization.

\begin{lstlisting}[style=ir]
define i32 @max(i32 %a, i32 %b) {
entry:
  %cmp = icmp sgt i32 %a, %b  ; Compare: a > b?
  br i1 %cmp, label %if.then, label %if.else

if.then:
  ret i32 %a

if.else:
  ret i32 %b
}
\end{lstlisting}

Each basic block has three essential properties: \textbf{\textit{unique label}}, \textbf{\textit{sequential instructions}}, and \textbf{\textit{terminator instruction}}. We examine each of these properties below.

\subsubsection{Unique Label}

Every basic block has a name (\texttt{entry}, \texttt{if.then}, \texttt{if.else}) that serves as a branch target. The label identifies the block and allows other instructions to reference it as a control flow destination. Labels must be unique within a function---if two labels with the same name appear in the same function, the compiler will reject the code with a redefinition error. Note that labels in different functions can share the same name, as label scope is local to each function. By convention, the first basic block is named \texttt{entry}.

Each instruction in a function must belong to a labeled basic block---there cannot be instructions floating outside of blocks. After a label, each instruction must appear on a new line; instructions cannot be placed on the same line as the label itself. While not syntactically required, instructions are typically indented below their block's label (usually with two spaces) to improve readability and visually distinguish labels from their instructions.

\begin{lstlisting}[style=ir]
entry:
  %x = add i32 1, 2        ; New line, indented (typical style)
  %y = mul i32 %x, 3
  br label %next

next:
%z = sub i32 %y, 1         ; New line, not indented (valid but unusual)
ret i32 %z
\end{lstlisting}

\subsubsection{Instructions and Opcodes}

No control flow transfers occur in the middle of a basic block. Once execution enters a block at its label, all instructions execute sequentially from top to bottom until the terminator is reached. This property---single entry, single exit---is fundamental to control flow analysis and enables many compiler optimizations by simplifying data flow analysis.

\acr{LLVM} instructions are the fundamental operations that manipulate values. While the exact syntax varies by instruction type, most instructions share a common structure consisting of an \textbf{optional result variable}, an \textbf{opcode} (operation name), \textbf{type information}, and \textbf{operands}. The typical form for binary operations is:

\begin{lstlisting}[style=ir]
%result = opcode type operand1, operand2
\end{lstlisting}

For example, \texttt{\%sum = add i32 \%a, \%b} follows this pattern. However, different instruction families have variations: comparison instructions include a \textit{predicate} (\texttt{\%cmp = icmp eq i32 \%a, \%b}), memory operations use explicit pointer types (\texttt{\%val = load i32, ptr \%ptr}), and side-effecting instructions like \texttt{store} produce no result. Despite these variations, all instructions combine an operation identifier with typed operands to perform their specific computation or effect.

The complete and authoritative documentation for all \acr{LLVM} \acr{IR} instructions, including detailed syntax, semantics, and behavior specifications, is provided in the \textbf{\acr{LLVM} Language Reference Manual}, the official documentation maintained by the \acr{LLVM} Project\footnote{Available online at \url{https://llvm.org/docs/LangRef.html}}. This comprehensive reference covers every instruction opcode, from basic arithmetic operations to advanced vector manipulations and atomic operations. The canonical instruction enumeration is maintained in the \acr{LLVM} Project source code in the file \texttt{Instruction.def}\footnote{In the \acr{LLVM} source tree, this is located at \texttt{llvm/include/llvm/IR/Instruction.def}. On Ubuntu systems with \acr{LLVM} installed, you can find it at \texttt{/usr/include/llvm-<version>/llvm/IR/Instruction.def} (e.g., \texttt{/usr/include/llvm-18/llvm/IR/Instruction.def}). This file contains C++ macro definitions that enumerate all instruction opcodes. You can count all instructions on Linux with: \texttt{grep "HANDLE.*INST" /usr/include/llvm-18/llvm/IR/Instruction.def | wc -l}}. Note that unlike assembly languages for specific architectures, \acr{LLVM} does not provide a command-line tool to enumerate all opcodes---the Language Reference Manual serves as the definitive catalog. What follows is a concise overview of the most commonly used instruction categories with representative examples to illustrate typical usage patterns.

\textbf{Arithmetic instructions:}

\begin{lstlisting}[style=ir]
%sum = add i32 %a, %b         ; Integer addition
%diff = sub i32 %a, %b        ; Integer subtraction
%prod = mul i32 %a, %b        ; Integer multiplication
%quot = sdiv i32 %a, %b       ; Signed division
%rem = srem i32 %a, %b        ; Signed remainder

%fsum = fadd float %x, %y     ; 32-bit float addition
%fdiff = fsub float %x, %y    ; 32-bit float subtraction
%fprod = fmul float %x, %y    ; 32-bit float multiplication
%fquot = fdiv float %x, %y    ; 32-bit float division

%dsum = fadd double %x, %y    ; 64-bit double addition
%ddiff = fsub double %x, %y   ; 64-bit double subtraction
%dprod = fmul double %x, %y   ; 64-bit double multiplication
%dquot = fdiv double %x, %y   ; 64-bit double division
\end{lstlisting}

Note that floating-point opcodes (\texttt{fadd}, \texttt{fsub}, \texttt{fmul}, \texttt{fdiv}) work for all floating-point types---the type specifier (\texttt{float}, \texttt{double}, \texttt{half}, \texttt{fp128}) determines the precision.

\textbf{Advanced arithmetic operations:}

For transcendental functions and advanced mathematical operations, \acr{LLVM} provides intrinsic functions. These are built-in functions prefixed with \texttt{llvm.} that map to efficient hardware instructions or optimized library implementations.

\begin{lstlisting}[style=ir]
; Trigonometric functions (arguments in radians)
%s = call float @llvm.sin.f32(float %angle)
%c = call double @llvm.cos.f64(double %angle)
%t = call float @llvm.tan.f32(float %x)

; Inverse trigonometric functions
%asin = call float @llvm.asin.f32(float %x)
%acos = call double @llvm.acos.f64(double %x)
%atan = call float @llvm.atan.f32(float %x)
%atan2 = call double @llvm.atan2.f64(double %y, double %x)

; Hyperbolic functions
%sinh = call float @llvm.sinh.f32(float %x)
%cosh = call double @llvm.cosh.f64(double %x)
%tanh = call float @llvm.tanh.f32(float %x)

; Exponential and logarithmic functions
%exp = call float @llvm.exp.f32(float %x)       ; e^x
%exp2 = call double @llvm.exp2.f64(double %x)   ; 2^x
%exp10 = call float @llvm.exp10.f32(float %x)   ; 10^x
%log = call double @llvm.log.f64(double %x)     ; natural log (ln)
%log2 = call float @llvm.log2.f32(float %x)     ; log base 2
%log10 = call double @llvm.log10.f64(double %x) ; log base 10

; Power and root functions
%sqrt = call float @llvm.sqrt.f32(float %x)
%pow = call double @llvm.pow.f64(double %base, double %exp)
%powi = call float @llvm.powi.f32(float %base, i32 %exp)

; Rounding and absolute value
%floor = call float @llvm.floor.f32(float %x)
%ceil = call double @llvm.ceil.f64(double %x)
%round = call float @llvm.round.f32(float %x)
%abs = call double @llvm.fabs.f64(double %x)

; Fused multiply-add (a * b + c in one operation)
%fma = call float @llvm.fma.f32(float %a, float %b, float %c)
\end{lstlisting}

These intrinsics follow the naming convention \texttt{@llvm.<operation>.<type>}, where the type suffix indicates the floating-point precision (\texttt{f32} for \texttt{float}, \texttt{f64} for \texttt{double}, \texttt{f16} for \texttt{half}). Note that all trigonometric functions operate in radians, not degrees.

For mathematical functions not provided as \acr{LLVM} intrinsics, you must declare and call external library functions from the C math library. Notably, while hyperbolic functions (\texttt{sinh}, \texttt{cosh}, \texttt{tanh}) are available as intrinsics, their inverses (\texttt{asinh}, \texttt{acosh}, \texttt{atanh}) are not:

\begin{lstlisting}[style=ir]
; Declare external inverse hyperbolic functions
declare float @asinhf(float)
declare double @asinh(double)
declare float @acoshf(float)
declare double @acosh(double)
declare float @atanhf(float)
declare double @atanh(double)

; Call the external functions
%result_f = call float @asinhf(float %x)
%result_d = call double @asinh(double %x)
%acosh_result = call double @acosh(double %x)
%atanh_result = call float @atanhf(float %x)
\end{lstlisting}

The \texttt{declare} statement tells \acr{LLVM} that these functions exist externally, but they are not implemented in your \acr{IR} code. At link time, you must link against the C math library (\texttt{libm}) to resolve these function references. For example, when compiling to an executable:

\begin{lstlisting}[style=ir]
# Compile LLVM IR to object file
llc -filetype=obj myfile.ll -o myfile.o

# Link with math library (-lm flag)
clang myfile.o -lm -o myprogram
\end{lstlisting}

Or directly with \texttt{clang}:

\begin{lstlisting}[style=ir]
clang myfile.ll -lm -o myprogram
\end{lstlisting}

The \texttt{-lm} flag instructs the linker to include the math library where the implementations of \texttt{asinh}, \texttt{acosh}, \texttt{atanh}, and other standard math functions are found.

\textbf{Bitwise operations:}

\begin{lstlisting}[style=ir]
%and_result = and i32 %a, %b  ; Bitwise AND
%or_result = or i32 %a, %b    ; Bitwise OR
%xor_result = xor i32 %a, %b  ; Bitwise XOR
%shift_left = shl i32 %a, 2   ; Shift left by 2 bits
%shift_right = lshr i32 %a, 2 ; Logical shift right
%arith_shift = ashr i32 %a, 2 ; Arithmetic shift right (sign-extend)
\end{lstlisting}

\textbf{Comparison instructions:}

\begin{lstlisting}[style=ir]
; Integer comparisons (icmp)
%eq = icmp eq i32 %a, %b      ; Equal
%ne = icmp ne i32 %a, %b      ; Not equal
%gt = icmp sgt i32 %a, %b     ; Signed greater than
%lt = icmp slt i32 %a, %b     ; Signed less than
%uge = icmp uge i32 %a, %b    ; Unsigned greater or equal

; Floating-point comparisons (fcmp)
%feq = fcmp oeq float %x, %y  ; Ordered equal
%fgt = fcmp ogt float %x, %y  ; Ordered greater than
%fne = fcmp une float %x, %y  ; Unordered or not equal
\end{lstlisting}

\textbf{Memory operations:}

\begin{lstlisting}[style=ir]
%ptr = alloca i32             ; Stack allocation
store i32 42, ptr %ptr        ; Write to memory
%val = load i32, ptr %ptr     ; Read from memory

; Pointer arithmetic with getelementptr (GEP)
%array_ptr = getelementptr [10 x i32], ptr %base, i64 0, i64 3
; Access element 3 of array at %base
\end{lstlisting}

\textbf{Type conversion instructions:}

\begin{lstlisting}[style=ir]
%ext = zext i32 %a to i64     ; Zero-extend to wider type
%trunc = trunc i64 %b to i32  ; Truncate to narrower type
%fp_to_int = fptosi float %x to i32  ; Float to signed int
%int_to_fp = sitofp i32 %y to float  ; Signed int to float
%bitcast = bitcast i32 %z to float   ; Reinterpret bits
\end{lstlisting}

\textbf{Function calls:}

\begin{lstlisting}[style=ir]
%result = call i32 @compute(i32 %x, i32 %y)
call void @print_value(i32 %result)
\end{lstlisting}

\textbf{Vector operations (\acr{SIMD}):}

\begin{lstlisting}[style=ir]
%vec_a = <4 x float> <float 1.0, float 2.0, float 3.0, float 4.0>
%vec_b = <4 x float> <float 5.0, float 6.0, float 7.0, float 8.0>
%vec_sum = fadd <4 x float> %vec_a, %vec_b
\end{lstlisting}

\subsubsection{Terminator Instructions}

Every basic block must end with exactly one terminator that transfers control to another block or exits the function. The terminator determines where execution continues after the current block completes. The complete list of terminator instructions is:

\begin{itemize}[noitemsep]
  \item \texttt{ret <type> <value>} --- Return from function with a value
  \item \texttt{ret void} --- Return from function without a value
  \item \texttt{br label \%dest} --- Unconditional jump to a basic block
  \item \texttt{br i1 \%cond, label \%true, label \%false} --- Conditional branch based on boolean condition
  \item \texttt{switch <type> <value>, label <default> [ <cases> ]} --- Multi-way branch (jump table) that dispatches based on integer value
  \item \texttt{indirectbr <type> <address>, [ <labels> ]} --- Indirect branch through a computed address (used for computed gotos and jump tables)
  \item \texttt{invoke <return\_type> <function> (<args>) to label \%normal unwind label \%exception} --- Function call with exception handling support; transfers to \texttt{\%normal} on success or \texttt{\%exception} on exception
  \item \texttt{resume <type> <value>} --- Resume exception propagation (used in exception handling to re-throw)
  \item \texttt{catchswitch within <parent> [ <handlers> ] unwind <default>} --- Exception catch dispatch that transfers control to appropriate catch handler
  \item \texttt{catchret from <catchpad> to label <successor>} --- Return from catch handler to normal control flow
  \item \texttt{cleanupret from <cleanuppad> unwind <destination>} --- Return from cleanup handler (finally block) to either normal flow or exception propagation
  \item \texttt{unreachable} --- Marks code paths that should never be reached; enables optimizations by indicating dead code
  \item \texttt{callbr <return\_type> <function> (<args>) to label \%fallthrough [<indirect\_labels>]} --- Inline assembly call with multiple possible successors (supports \texttt{asm goto})
\end{itemize}

The first basic block in every function is conventionally named \texttt{entry} and serves as the function's entry point.

\subsection{The \acr{PHI} Instruction}

When control flow merges from multiple basic blocks, \acr{SSA} form requires a mechanism to select the appropriate value based on which predecessor block was executed. The \texttt{phi} instruction serves this purpose.

\begin{lstlisting}[style=ir]
define i32 @abs(i32 %x) {
entry:
  %is_neg = icmp slt i32 %x, 0
  br i1 %is_neg, label %negate, label %done

negate:
  %neg = sub i32 0, %x
  br label %done

done:
  %result = phi i32 [ %x, %entry ], [ %neg, %negate ]
  ret i32 %result
}
\end{lstlisting}

The \texttt{phi} instruction syntax is:

\begin{lstlisting}[style=ir]
%result = phi <type> [ <value1>, %block1 ], [ <value2>, %block2 ], ...
\end{lstlisting}

Semantics: "If control flow arrived from \texttt{\%block1}, use \texttt{<value1>}; if from \texttt{\%block2}, use \texttt{<value2>}."

\textbf{Why \texttt{phi} is necessary:}

Without \texttt{phi}, we cannot maintain \acr{SSA} form when values merge from different control paths. Consider this invalid attempt:

\begin{lstlisting}[style=ir]
; INVALID: Cannot assign to %result in multiple blocks
entry:
  %result = ... ; First assignment
  br ...

other_block:
  %result = ... ; ERROR: violates SSA (redefinition)
\end{lstlisting}

The \texttt{phi} instruction solves this by creating a single assignment point that selects from multiple incoming values.

\textbf{Multiple \texttt{phi} nodes:}

A block can have multiple \texttt{phi} instructions. All \texttt{phi} instructions in a block are conceptually executed simultaneously before any other instructions:

\begin{lstlisting}[style=ir]
loop:
  %i = phi i32 [ 0, %entry ], [ %i_next, %loop ]
  %sum = phi i32 [ 0, %entry ], [ %sum_next, %loop ]
  %i_next = add i32 %i, 1
  %sum_next = add i32 %sum, %i
  %cond = icmp slt i32 %i_next, 10
  br i1 %cond, label %loop, label %exit
\end{lstlisting}

This example implements a loop that sums integers from 0 to 9, demonstrating how \texttt{phi} enables iterative control flow in \acr{SSA} form.
